\documentclass[]{article}

\usepackage{graphicx} 			% Imagens
\usepackage[utf8]{inputenc}		% Acentos
\usepackage{subfigure}			% Conjuntos de figuras
\usepackage{tikz}				% Desenhar no Latex

%opening
\title{Trabalhando com Imagens}
\author{Lucas Arantes Berg}

\begin{document}

\maketitle

\begin{abstract}
	Aqui vai o resumo.
\end{abstract}

\section{Teste de imagens}

% h = here
% b = bottom of the page
% p = separate page
% t = top of page
% ! = force position
\begin{figure}[hb]
	\includegraphics[width=2cm]{figs/smiley.jpeg}
	\caption{Sorria você não está sendo filmado.}
	\label{fig:smiley}
\end{figure}

\begin{figure}[hb]
	\centering
	\subfigure[ref1][Feliz]{\includegraphics[width=2cm]{figs/smiley.jpeg}}
	\qquad
	\subfigure[ref1][Triste]{\includegraphics[width=2cm]{figs/sad-smiley.png}}
	\caption{Smiles.}
	\label{fig:smiles}
\end{figure}

\begin{figure}[hb]
	\centering
	\begin{tikzpicture}
		\draw (0,0) -- (0,1) -- (1,0) -- cycle;
	\end{tikzpicture}
	\caption{Um triângulo louco.}
	\label{fig:triangulo}	
\end{figure}

\begin{figure}[hb]
	\centering
	\begin{tikzpicture}[rotate=90]
	\draw (0:1) -- (2*72:1) -- (4*72:1) -- (6*72:1) -- (8*72:1) -- cycle;
	\end{tikzpicture}
	\caption{Um estrela maneira.}
	\label{fig:estrela}	
\end{figure}

\newpage
\begin{tabular}{|c|r|l|}
	\hline
	a & bb & ccc \\ \hline
	bb & ccc & a \\ \hline
	ccc & a & bb \\ \hline
\end{tabular}




\end{document}
